\documentclass[a4paper,12pt,oneside]{article}
\makeatletter
\def\input@path{{../../}}
\makeatother
\usepackage{discrete-practice}

\begin{document}

\PracticeHeader{1}{7}{Карты Карно и минимизация}

\Task{Задание 1}
Нахождение СДНФ и СКНФ при помощи карт Карно.
\begin{numlist}
    \item $f=A\overline{B}+B\overline{D}+\overline{A}C$
    \item $f=BD+AC+CD$
    \item $f=\overline{D}+\overline{A}C+\overline{B}D$
    \item $f=\overline{C}+BD+\overline{A}\,\overline{B}$
    \item $f=ABC+BD$
    \item $f=A\overline{C}+AC+BD$
\end{numlist}

\Task{Задание 2}
Нанесение функции на карту Карно:
\[
f=\overline{C}+\overline{A}B+A\overline{B},\qquad
f=\overline{A}B+BCD+\overline{A}\,\overline{C}
\]

\Task{Задание 3}
Алгебраическое упрощение:
\[
f=A\overline{C}+BC+\overline{A}B
\]

\Task{Задание 4}
Импликанты. Чему равна длина вектора минтермов функции, если число её импликант равно:
\[
1)\ 512 \qquad 2)\ 8 \qquad 3)\ 1
\]

\Task{Задание 5}
Методом Квайна найдите сокращённую ДНФ. Укажите число простых импликант, из которых состоит сокращённая форма, и число вхождений переменных:
\[
f(A,B,C)=(0,3,5,6,7)
\]
Что делать, если необходимо найти сокращённую КНФ?

\Task{Задание 6}
Функция задана списком номеров минтермов. Найти минимальную ДНФ и КНФ функции:
\[
f(A,B,C,D)=(0,1,2,6,7,8,12,15)
\]
Может ли у функции не быть ядерных импликант?
\begin{center}
  \raisebox{-.5\height}{\includegraphics[width=.55\linewidth]{P7.image1.png}}\hfill
  \raisebox{-.5\height}{\includegraphics[width=.18\linewidth]{P7.image2.png}}\hfill
  \raisebox{-.5\height}{\includegraphics[width=.18\linewidth]{P7.image3.png}}
\end{center}

\end{document}
